\chapter[\emj{🎒}~DP: Knapsack, Subset Sum y Coin Change]{\emj{🎒}~DP: Knapsack, Subset Sum y Coin Change}

\begin{dsaquees}

Una mochila 🎒 que aguanta cierto peso máximo y un montón de objetos, cada
uno con su peso y su valor. ¿Qué combinación METES para maximizar el
valor sin pasarte del peso?

No puedes ser goloso (agarrar lo más valioso) porque a veces dos objetos
medianos ganan a uno grande. Hay que PROBAR combinaciones, pero de forma
inteligente: recordando la mejor respuesta para cada "peso disponible" y
reusándola. Eso es Programación Dinámica: no recalcular lo ya resuelto.

\end{dsaquees}

\begin{dsanino}

Tienes una mochila que solo aguanta 5 kilos y varios juguetes 🧸. Cada
juguete pesa algo y te gusta más o menos (su "valor"). Quieres llevarte
los juguetes que MÁS te gusten sin que la mochila pese más de 5 kilos.

El truco es hacer una tabla. Vas juguete por juguete preguntando: "si
tuviera espacio para 1 kilo, para 2, para 3... ¿qué es lo mejor que puedo
llevar?". Para cada juguete decides: ¿lo dejo fuera, o lo meto y sumo su
valor? Te quedas con la opción que dé MÁS valor. Como escribes las
respuestas en la tabla, nunca tienes que repetir la misma cuenta dos
veces. Al final, la esquina de la tabla te dice el mejor botín posible.

\end{dsanino}

\begin{dsaotraforma}

El \textbf{0/1 Knapsack} es el problema base de una familia enorme de DP.
Con $n$ objetos (peso $w_i$, valor $v_i$) y capacidad $W$, maximiza el
valor total eligiendo un subconjunto cuyo peso no supere $W$. "0/1"
porque cada objeto se toma entero o no se toma.

La recurrencia, con \verb|dp[i][c]| = mejor valor usando los primeros $i$
objetos con capacidad $c$:
\[ dp[i][c] = \max\big(dp[i-1][c],\ \ v_i + dp[i-1][c - w_i]\big) \]
(no tomar el objeto $i$, o tomarlo si cabe).

Variantes que son el MISMO esqueleto:
\begin{itemize}
\item \textbf{Subset Sum / Partition}: ¿hay un subconjunto que sume $S$?
      (knapsack booleano).
\item \textbf{Coin Change}: mínimo de monedas para un monto (unbounded
      knapsack: cada moneda se repite).
\end{itemize}

\end{dsaotraforma}

\begin{dsadiagrama}
\begin{verbatim}
0/1 KNAPSACK  (W = 5)     objetos: (peso, valor)
  o1=(2,3)  o2=(3,4)  o3=(4,5)  o4=(5,6)

Tabla dp[objeto][capacidad], mejor valor:

 cap:  0  1  2  3  4  5
  {}   0  0  0  0  0  0
  o1   0  0  3  3  3  3
  o2   0  0  3  4  4  7   <- o1+o2 = 3+4 = 7 con peso 5
  o3   0  0  3  4  5  7
  o4   0  0  3  4  5  7

Respuesta = dp[n][W] = 7  (llevar o1 y o2)
Cada celda: max(de arriba, valor + celda arriba-desplazada por su peso)
\end{verbatim}
\end{dsadiagrama}

\begin{lstlisting}
CÓMO FUNCIONA (0/1 knapsack)
- dp[c] = mejor valor con capacidad c (versión 1D optimizada).
- Por cada objeto, recorre la capacidad de W hacia ABAJO (para no reusar
  el mismo objeto dos veces).
- dp[c] = max(dp[c], valor + dp[c - peso]).

PSEUDOCÓDIGO (0/1 knapsack 1D)
──────────────────────────────────────────────────────────────

función knapsack(pesos, valores, W):
    dp = arreglo de ceros de tamaño W+1
    PARA i desde 0 hasta n-1:
        PARA c desde W hasta pesos[i]:      // ¡hacia abajo!
            dp[c] = max(dp[c], valores[i] + dp[c - pesos[i]])
    return dp[W]

CÓDIGO C++ (0/1 knapsack, 1D)
──────────────────────────────────────────────────────────────

#include <vector>
#include <algorithm>
using namespace std;

int knapsack(vector<int>& peso, vector<int>& valor, int W) {
    vector<int> dp(W + 1, 0);
    for (int i = 0; i < (int)peso.size(); i++)
        for (int c = W; c >= peso[i]; c--)          // hacia abajo: 0/1
            dp[c] = max(dp[c], valor[i] + dp[c - peso[i]]);
    return dp[W];
}

CÓDIGO C++ (Subset Sum: ¿existe subconjunto que sume S?)
──────────────────────────────────────────────────────────────

bool subsetSum(vector<int>& nums, int S) {
    vector<bool> dp(S + 1, false);
    dp[0] = true;                          // suma 0 siempre posible
    for (int x : nums)
        for (int c = S; c >= x; c--)       // hacia abajo
            dp[c] = dp[c] || dp[c - x];
    return dp[S];
}

CÓDIGO C++ (Coin Change: mínimo de monedas, unbounded)
──────────────────────────────────────────────────────────────

int coinChange(vector<int>& coins, int amount) {
    vector<int> dp(amount + 1, amount + 1);      // "infinito"
    dp[0] = 0;
    for (int c : coins)
        for (int a = c; a <= amount; a++)        // hacia ARRIBA: repetible
            dp[a] = min(dp[a], dp[a - c] + 1);
    return dp[amount] > amount ? -1 : dp[amount];
}

COMPLEJIDAD (Big-O)
──────────────────────────────────────────────────────────────

  Problema      │ Tiempo     │ Espacio
  ──────────────┼────────────┼──────────
  0/1 Knapsack  │ O(n·W)     │ O(W)
  Subset Sum    │ O(n·S)     │ O(S)
  Coin Change   │ O(n·A)     │ O(A)

* Es "pseudo-polinómico": depende del VALOR W, no solo de n.
\end{lstlisting}

\begin{dsaentrevistas}

✅ Dirección del bucle interno = la clave: hacia ABAJO (W..peso) para
   0/1 (cada objeto una vez); hacia ARRIBA (peso..W) para unbounded
   (objeto repetible, como monedas). Confundirlas cambia el resultado.

💡 "Partition Equal Subset Sum" (LC 416): ¿se puede partir en dos mitades
   de igual suma? Es Subset Sum con $S = \mathrm{total}/2$. Si el total es
   impar, imposible de una.

⚠️ Pseudo-polinómico: O(n·W) parece polinomial pero W puede ser enorme.
   Si W es gigante y n pequeño, quizá otra técnica convenga.

🔑 Reconstruir la solución: si te piden CUÁLES objetos, guarda la tabla
   2D y recórrela hacia atrás comparando dp[i][c] con dp[i-1][c].

\end{dsaentrevistas}

\begin{dsaresumen}

• 0/1 Knapsack: max valor sin pasar la capacidad; cada objeto entero o no.
• Recurrencia: max(no tomar, valor + tomar si cabe).
• Optimización 1D: bucle de capacidad hacia ABAJO para 0/1.
• Subset Sum / Partition: knapsack booleano (¿suma S posible?).
• Coin Change: unbounded knapsack; bucle hacia ARRIBA (repetible).
• Complejidad pseudo-polinómica O(n·W).
════════════════════════════════════════════════════════════════

\end{dsaresumen}
